\section{Anforderungen und Qualitätsziele}
\label{chap:requirements-quality-goals}

Die Qualitätsziele sind Portabilität, Produktschutz, Idempotenz, Rollback,
Plattformunabhängigkeit, reproduzierbarer Export, versionierte Migration und
überprüfbare Dokumentation. Ein \CheckMode{} muss lesend bleiben; ein \FullFix{} darf
nur eine zuvor sichtbar geplante Operation ausführen.

\CaseTable{
\begin{tabularx}{0.94\textwidth}{@{}lX@{}}
\toprule
Ziel & Prüfkriterium \\
\midrule
Portabilität & Alle Laufzeitpfade stammen aus \ProjectContext{}. \\
Produktschutz & Unbekannte und produktverwaltete Dateien bleiben unverändert. \\
Idempotenz & Ein zweiter erfolgreicher \FullFix{} plant keine Operationen. \\
Rollback & Fehlerhafte Veröffentlichung stellt erfasste Vorzustände wieder her. \\
\bottomrule
\end{tabularx}
}{Qualitätsziele und beobachtbare Prüfkriterien.}{tab:quality-goals}

Diese Ziele beschreiben Anforderungen, keine pauschale Behauptung über jede künftige
Adapterimplementierung. Die zugehörige Strukturprüfung ist bis zum nächsten
eindeutigen Commit als \Planned{} (\evidence{EV-DOCS-001}) registriert.
